\documentclass[11pt,a4paper]{article}
\usepackage{fontspec,amsmath,amssymb,amsthm,microtype}
\setmainfont{lmroman10-regular.otf}[BoldFont=lmroman10-bold.otf,ItalicFont=lmroman10-italic.otf,BoldItalicFont=lmroman10-bolditalic.otf]
\usepackage{xcolor}
\usepackage[most]{tcolorbox}
\usepackage{mathrsfs}
\providecommand{\lt}{<}
\providecommand{\gt}{>}
\setmonofont{DejaVuSansMono.ttf}[BoldFont=DejaVuSansMono-Bold.ttf,ItalicFont=DejaVuSansMono-Oblique.ttf,BoldItalicFont=DejaVuSansMono-BoldOblique.ttf,Scale=0.86]
\newfontfamily\leanfallback{STIXTwoMath-Regular.otf}[Scale=0.86]
\definecolor{leanrule}{gray}{0.62}
\definecolor{leanback}{gray}{0.962}
\definecolor{leantitle}{gray}{0.22}
\definecolor{leanmuted}{gray}{0.38}
\newtcolorbox{leanbox}[1]{enhanced,breakable,colback=leanback,colframe=leanrule,
  boxrule=0.4pt,arc=0.6mm,left=2.4mm,right=2.4mm,top=1.8mm,bottom=1.8mm,
  before skip=0.8em,after skip=1em,fonttitle=\small\bfseries,coltitle=leantitle,
  colbacktitle=leanback,titlerule=0.3pt,title={#1}}
\newlength{\leancharwidth}
\newenvironment{leancode}{\par\smallskip\begingroup\ttfamily\footnotesize\raggedright
  \settowidth{\leancharwidth}{M}\setlength{\parskip}{0pt}\setlength{\parindent}{0pt}
  \frenchspacing}{\par\endgroup\smallskip}
\newcommand{\leanline}[2]{\hangindent=\dimexpr#1\leancharwidth+4\leancharwidth\relax
  \hangafter=1\noindent\hspace*{#1\leancharwidth}#2\par}
\newcommand{\leansource}[1]{\par{\footnotesize\color{leanmuted}#1}\par}
\newcommand{\leanlabel}[1]{\par\smallskip\noindent{\small\bfseries #1}\par}
\usepackage[margin=25mm]{geometry}
\usepackage[unicode,hidelinks]{hyperref}
\hypersetup{pdfauthor={LangevinLean Project}}
\setlength{\emergencystretch}{3em}
\setlength{\parskip}{0.4em}
\setlength{\parindent}{0pt}
\allowdisplaybreaks[2]
\raggedbottom
\title{\texorpdfstring{The shifted energy floor and the coefficient six}{The shifted energy floor and the coefficient six}}
\author{LangevinLean Project}
\date{}
\begin{document}\maketitle
Draft mathematical exposition; these notes may contain mistakes.


Let \(q=1-\alpha h\), \(C=h^2\beta d\), \(F=19h^3\beta^2d\) and \(E_*=19h^2\kappa\beta d\), with \(\kappa=\beta/\alpha\). The exact fixed-point identity \(qE_*+F=E_*\) turns a supplied scalar energy recurrence into uniform and finite-time bounds. The final norm comparison uses only triangle inequalities and these two squared budgets.



Earlier results used here: \cite[Section~3.15]{LangevinLean2026ULASC}.

\section{\texorpdfstring{The budgets and their exact fixed point}{The budgets and their exact fixed point}}


\textbf{Claim.} Under the numerical parameter restrictions, \(C\ge0\), \(E_*\ge0\), \(C\le E_*\), and \(qE_*+F=E_*\).




\textbf{Proof.} The budgets are products of nonnegative factors. Since \(\kappa\ge1\), \(E_*=19\kappa C\ge C\). For the fixed point, \((1-q)E_*=\alpha h\cdot19h^2(\beta/\alpha)\beta d=19h^3\beta^2d=F\). Division by \(\alpha\) is valid because \(\alpha>0\).

\begin{leanbox}{Lean: \texttt{Parameters.\allowbreak{}correctorBudget}, \texttt{Parameters.\allowbreak{}energyForcing}, \texttt{Parameters.\allowbreak{}energyFloor} and 4 more}
\leansource{\texttt{Parameters.\allowbreak{}correctorBudget} \quad WassersteinMixing/\allowbreak{}Recurrence.lean, lines 17–18}
\begin{leancode}
\leanline{0}{def\ Parameters.correctorBudget\ (P\ :\ Parameters)\ :\ ℝ\ :=}
\leanline{2}{P.step\ \textasciicircum{}\ 2\ *\ P.beta\ *\ P.dimension}
\end{leancode}
\leansource{\texttt{Parameters.\allowbreak{}energyForcing} \quad WassersteinMixing/\allowbreak{}Recurrence.lean, lines 21–22}
\begin{leancode}
\leanline{0}{def\ Parameters.energyForcing\ (P\ :\ Parameters)\ :\ ℝ\ :=}
\leanline{2}{19\ *\ P.step\ \textasciicircum{}\ 3\ *\ P.beta\ \textasciicircum{}\ 2\ *\ P.dimension}
\end{leancode}
\leansource{\texttt{Parameters.\allowbreak{}energyFloor} \quad WassersteinMixing/\allowbreak{}Recurrence.lean, lines 25–26}
\begin{leancode}
\leanline{0}{def\ Parameters.energyFloor\ (P\ :\ Parameters)\ :\ ℝ\ :=}
\leanline{2}{19\ *\ P.step\ \textasciicircum{}\ 2\ *\ P.condition\ *\ P.beta\ *\ P.dimension}
\end{leancode}
\leansource{\texttt{Parameters.\allowbreak{}correctorBudget\_\allowbreak{}nonneg} \quad WassersteinMixing/\allowbreak{}Recurrence.lean, lines 28–30}
\begin{leancode}
\leanline{0}{theorem\ Parameters.correctorBudget\_nonneg\ (P\ :\ Parameters)\ :\ 0\ ≤\ P.correctorBudget}
\end{leancode}
\leansource{\texttt{Parameters.\allowbreak{}energyFloor\_\allowbreak{}nonneg} \quad WassersteinMixing/\allowbreak{}Recurrence.lean, lines 32–36}
\begin{leancode}
\leanline{0}{theorem\ Parameters.energyFloor\_nonneg\ (P\ :\ Parameters)\ :\ 0\ ≤\ P.energyFloor}
\end{leancode}
\leansource{\texttt{Parameters.\allowbreak{}energy\_\allowbreak{}fixed\_\allowbreak{}point} \quad WassersteinMixing/\allowbreak{}Recurrence.lean, lines 39–44}
\begin{leancode}
\leanline{0}{theorem\ Parameters.energy\_fixed\_point\ (P\ :\ Parameters)\ :}
\leanline{4}{P.contraction\ *\ P.energyFloor\ +\ P.energyForcing\ =\ P.energyFloor}
\end{leancode}
\leansource{\texttt{Parameters.\allowbreak{}correctorBudget\_\allowbreak{}le\_\allowbreak{}floor} \quad WassersteinMixing/\allowbreak{}Recurrence.lean, lines 46–56}
\begin{leancode}
\leanline{0}{theorem\ Parameters.correctorBudget\_le\_floor\ (P\ :\ Parameters)\ :}
\leanline{4}{P.correctorBudget\ ≤\ P.energyFloor}
\end{leancode}
\leanlabel{How the Lean proof follows the argument}
The three parameter definitions retain the exact constants. The sign lemmas use positive curvature and nonnegative dimension. \texttt{energy\_\allowbreak{}fixed\_\allowbreak{}point} clears the nonzero denominator and proves the polynomial identity. \texttt{corrector\allowbreak{}Budget\_\allowbreak{}le\_\allowbreak{}floor} rewrites the floor as \(19\kappa\) times the budget.
\end{leanbox}

\section{\texorpdfstring{Uniform and finite-time energy envelopes}{Uniform and finite-time energy envelopes}}


\textbf{Claim.} Suppose \(E_{n+1}\le qE_n+F\). If \(E_0\le C\), then \(E_n\le E_*\) for every \(n\). Without the initial restriction,

\[E_n\le q^nE_0+(1-q^n)E_*.\]


\textbf{Proof.} The uniform bound follows by induction. Initially \(E_0\le C\le E_*\). Multiplication of the induction hypothesis by \(q\ge0\), followed by the recurrence and fixed-point identity, gives \(E_{n+1}\le qE_*+F=E_*\). For the finite-time bound, iterating the affine recurrence yields its geometric sum. Substitute \(F=(1-q)E_*\); since \(q<1\), cancellation of \(1-q\) gives the displayed expression.

\begin{leanbox}{Lean: \texttt{shifted\_\allowbreak{}energy\_\allowbreak{}uniform}, \texttt{shifted\_\allowbreak{}energy\_\allowbreak{}finite}}
\leansource{\texttt{shifted\_\allowbreak{}energy\_\allowbreak{}uniform} \quad WassersteinMixing/\allowbreak{}Recurrence.lean, lines 60–69}
\begin{leancode}
\leanline{0}{theorem\ shifted\_energy\_uniform\ (P\ :\ Parameters)\ (energy\ :\ ℕ\ →\ ℝ)}
\leanline{4}{(hzero\ :\ energy\ 0\ ≤\ P.correctorBudget)}
\leanline{4}{(hstep\ :\ ∀\ n,\ energy\ (n\ +\ 1)\ ≤\ P.contraction\ *\ energy\ n\ +\ P.energyForcing)}
\leanline{4}{(n\ :\ ℕ)\ :\ energy\ n\ ≤\ P.energyFloor}
\end{leancode}
\leansource{\texttt{shifted\_\allowbreak{}energy\_\allowbreak{}finite} \quad WassersteinMixing/\allowbreak{}Recurrence.lean, lines 72–84}
\begin{leancode}
\leanline{0}{theorem\ shifted\_energy\_finite\ (P\ :\ Parameters)\ (energy\ :\ ℕ\ →\ ℝ)}
\leanline{4}{(hstep\ :\ ∀\ n,\ energy\ (n\ +\ 1)\ ≤\ P.contraction\ *\ energy\ n\ +\ P.energyForcing)}
\leanline{4}{(n\ :\ ℕ)\ :}
\leanline{4}{energy\ n\ ≤\ P.contraction\ \textasciicircum{}\ n\ *\ energy\ 0\ +}
\leanline{6}{(1\ -\ P.contraction\ \textasciicircum{}\ n)\ *\ P.energyFloor}
\end{leancode}
\leanlabel{How the Lean proof follows the argument}
\texttt{shifted\_\allowbreak{}energy\_\allowbreak{}uniform} is an explicit natural-number induction. \texttt{shifted\_\allowbreak{}energy\_\allowbreak{}finite} invokes the shared exact affine-recurrence formula and substitutes the proved fixed point. These are scalar implications. The stationary comparison derives their recurrence from the constructed corrector, actual martingale centering and genuine curvature bounds; the realized comparison supplies all primitive inputs from the actual potential dynamics.
\end{leanbox}


Earlier results used here: \cite[Section~3.15.2]{LangevinLean2026ULASC}.

\section{\texorpdfstring{Converting the two squared budgets to the literal bias constant}{Converting the two squared budgets to the literal bias constant}}


\textbf{Claim.} If \(\|z+g\|^2\le E_*\) and \(\|g\|^2\le C\), then

\[\|z\|\le6h\sqrt{\kappa\beta d}.\]


\textbf{Proof.} Put \(R=h\sqrt{\kappa\beta d}\ge0\). Then \(E_*=19R^2\), while \(C\le R^2\) because \(\kappa\ge1\). Hence \(\|z+g\|\le5R\), using \(19\le25\), and \(\|g\|\le R\). The triangle inequality for \(z=(z+g)-g\) gives \(\|z\|\le6R\). This proves the source's coefficient \(6\) directly, including a zero-dimensional budget.

\begin{leanbox}{Lean: \texttt{shifted\_\allowbreak{}norm\_\allowbreak{}le\_\allowbreak{}bias}}
\leansource{\texttt{shifted\_\allowbreak{}norm\_\allowbreak{}le\_\allowbreak{}bias} \quad WassersteinMixing/\allowbreak{}Recurrence.lean, lines 88–121}
\begin{leancode}
\leanline{0}{theorem\ shifted\_norm\_le\_bias\ \{H\ :\ Type*\}\ [NormedAddCommGroup\ H]}
\leanline{4}{(P\ :\ Parameters)\ (z\ g\ :\ H)}
\leanline{4}{(hu\ :\ ‖z\ +\ g‖\ \textasciicircum{}\ 2\ ≤\ P.energyFloor)}
\leanline{4}{(hg\ :\ ‖g‖\ \textasciicircum{}\ 2\ ≤\ P.correctorBudget)\ :\ ‖z‖\ ≤\ P.bias}
\end{leancode}
\leanlabel{How the Lean proof follows the argument}
\texttt{shifted\_\allowbreak{}norm\_\allowbreak{}le\_\allowbreak{}bias} proves the squared identities for the nonnegative radius, compares each squared norm using \texttt{sq\_\allowbreak{}le\_\allowbreak{}sq\ensuremath{{}_{0}}}, and applies \texttt{norm\_\allowbreak{}sub\_\allowbreak{}le}. The proof does not use an unproved manipulation of a limit superior.
\end{leanbox}

\begin{thebibliography}{ULA-LSI}
\bibitem[ULA-SC]{LangevinLean2026ULASC}
LangevinLean Project. \emph{Unadjusted Langevin under Strong Convexity}. LL-01, 2026. Draft report, content revision d34ead06987c. \url{https://odakyildiz.github.io/LangevinLean/docs/pdf/UnadjustedLangevin/FullProof.pdf}.
\end{thebibliography}

\end{document}
